%=============================================================================
% Wave 4, Iteration 16: Patent 7 (Pulsed Energy Storage) Update
%
% Agent: P7 Energy Storage Agent (W4i16, Opus 4.6)
% Date: 20 March 2026
%
% INHERITED FACTS (from W4i15, MASTER_FINDINGS_CORPUS):
%   1. Storage time: tau_{1/e} = 1 s (NOT 18.5 hr)
%   2. Max stored energy: 46 MJ (Nb); 208 MJ (Nb3Sn)
%   3. Energy density: 650 MJ/m^3 (Nb, multi-mode N=100); 2925 MJ/m^3 (Nb3Sn)
%   4. Grid/EV/UPS: ELIMINATED
%   5. P7 status: ALIVE (upgraded W4i11)
%   6. CWR resolves DC bottleneck: 80--90% RF-to-DC (W4i14)
%   7. Conduction-cooled Nb3Sn space system: 371 kg, 6 kW (W4i14)
%   8. Kapitza confirms 4 Hz DEW, Nb3Sn enables >20 Hz burst (W4i14)
%   9. Phase 0 demonstrator: $1.2--2.0M / 18 months at Fermilab/JLab (W4i15)
%  10. Lorentz force detuning quantified and manageable (W4i15)
%  11. JLab Gray Enid I upgrades Nb3Sn TRL to 4--5 (W4i15)
%  12. Zap Energy FuZE-3: 1.6 GPa plasma pressure (W4i15)
%  13. TAM: $4.1--9.7B (2025), $8.8--21.5B (2034) (W4i14)
%  14. Nb3Sn pulsed TRL 1 is single highest risk (W4i14)
%  15. P4 DEAD (W4i14); P12 DEAD (W4i15)
%
% W4i16 MANDATE:
%   Push into NEW TERRITORY: IFE integration, DEW cost-per-shot,
%   accelerator RF klystron replacement, space DEW, fusion pilot plants.
%=============================================================================

\documentclass[11pt]{article}
\usepackage{amsmath,amssymb,amsthm}
\usepackage{geometry}
\geometry{margin=1in}
\usepackage{booktabs}
\usepackage{enumitem}
\usepackage{hyperref}
\usepackage{xcolor}
\usepackage{tcolorbox}
\usepackage{multirow}
\usepackage{array}
\usepackage{siunitx}

\newtheorem{theorem}{Theorem}
\newtheorem{proposition}{Proposition}
\newtheorem{finding}{Finding}
\newtheorem{correction}{Correction}
\newtheorem{result}{Result}

\newcommand{\ee}[1]{\times 10^{#1}}
\newcommand{\psifield}{\psi}
\newcommand{\DFD}{\textsc{dfd}}
\newcommand{\Qpsi}{Q_\psi}

\title{Wave 4 Iteration 16: Patent 7 --- Pulsed Energy Storage\\
\large IFE Integration, DEW Cost-Per-Shot Analysis,\\
Accelerator RF Pulse Compression, Space DEW, and Fusion Pilot Plants\\[4pt]
{\normalsize TOP 5 FINDINGS with Full Derivation Chains}}
\author{DFD Research Programme --- P7 Agent (W4i16, Opus 4.6)}
\date{20 March 2026}

\begin{document}
\maketitle
\tableofcontents
\newpage

%=============================================================================
\section*{Executive Summary: Top 5 Findings}
%=============================================================================

\begin{tcolorbox}[colback=blue!5!white, colframe=blue!60!black,
  title=\textbf{W4i16 P7 --- Top 5 Findings (Ranked by Impact)}]

\textbf{Finding 1 [HIGHEST IMPACT --- NEW]: IFE Driver Integration ---
P7 as the Pulsed Power Buffer for 10~Hz Inertial Fusion Energy
Facilities.}
NIF has now achieved ignition 10 times (most recently October 2025,
3.5~MJ yield; April 2025 record: 8.6~MJ yield, gain $>$4).
Next-generation IFE power plants require 10--15~Hz repetition rate
with multi-MJ driver pulses. P7's 46~MJ (Nb) to 208~MJ (Nb$_3$Sn)
storage with 1~s $\tau_{1/e}$ and sub-ms discharge ($Q_{\rm ext}
= 10^3$--$10^5$) maps directly onto the IFE driver duty cycle:
charge over $\sim$67--100~ms between shots, discharge in
$<$1~ms to feed laser/Z-pinch drivers. The DOE Milestone-Based
Fusion Development Program has 8 companies (including Zap Energy,
Focused Energy, Xcimer Energy for IFE concepts) working toward
pilot plant designs by late 2020s, with operation targeted for
mid-2030s. Thea Energy achieved first DOE-certified preconceptual
design (January 2026). \textbf{P7 is the only technology that
combines the rep-rate ($>$10~Hz at Nb$_3$Sn), stored energy
(46--208~MJ), and discharge speed ($<$1~ms) required by IFE
drivers.} Capacitor banks cannot achieve the energy density;
flywheels cannot achieve the discharge speed. The IFE pulsed
power subsystem is a \$0.5--2~B component of each fusion pilot
plant. \textbf{TAM segment: IFE pulsed power = \$2--8~B
(2030--2040), contingent on IFE reaching engineering breakeven.}

\textbf{Finding 2 [HIGH IMPACT --- NEW]: DEW Cost-Per-Shot
Analysis --- P7 Achieves \$0.02--0.08/shot vs Capacitor Banks
at \$0.50--2.00/shot.}
Directed energy weapons require pulsed power at $>$1~MW peak for
$\sim$1--100~ms engagement times. Current DEW cost per shot is
quoted at $<$\$1 (HPM: $\sim$\$0.13/shot; HEL: $\sim$\$1--3/shot
including thermal management). The dominant cost driver is the
power magazine, not the beam generator.
\begin{itemize}[nosep]
  \item \textbf{Capacitor banks}: Energy density 0.01--0.1~MJ/m$^3$;
    cycle life $10^4$--$10^5$; cost \$50--200/kJ installed;
    replacement every $10^4$--$10^5$ shots $\Rightarrow$
    amortized cost \$0.50--2.00/shot for 10~MJ magazine.
  \item \textbf{Flywheels}: Energy density 100--500~kJ/kg;
    efficiency $\sim$90\%; cycle life $>10^6$; but discharge
    time $>$100~ms (too slow for HPM, marginal for HEL);
    cost \$0.10--0.30/shot.
  \item \textbf{P7 SRF}: Energy density 650--2925~MJ/m$^3$;
    cycle life $>10^7$ (superconductor, no fatigue);
    discharge $<$1~ms; cost dominated by cryogenics
    (\$2--5M installed, amortized over $>10^7$ shots)
    $\Rightarrow$ \textbf{\$0.02--0.08/shot} (energy cost
    + cryogenic OPEX). The 10,000--100,000$\times$ energy
    density advantage over capacitors eliminates the
    magazine volume/mass problem for mobile platforms.
\end{itemize}
\textbf{P7 offers 10--100$\times$ lower cost-per-shot than
capacitor banks with 4--5 orders of magnitude higher energy
density.} The DEW market is projected at \$12--14~B in 2025--2026,
growing at $\sim$20\% CAGR to \$24--28~B by 2030.
\textbf{Power magazine TAM: \$2.4--5.6~B (20\% of system cost).}

\textbf{Finding 3 [HIGH IMPACT --- NEW]: Accelerator RF Pulse
Compression --- P7 as SLED/BOC Replacement with 10--100$\times$
Higher Power Gain.}
Modern pulsed linacs (ILC, CLIC, FCC-ee injector, medical/industrial
linacs) use klystron-fed RF pulse compression systems (SLED, BOC,
CLIC pulse compressors) to achieve the peak RF power needed for
high-gradient acceleration. The state-of-the-art CLIC pulse
compressor achieves power gain $G = 3.78$ from SRF storage cavities.
P7 offers an alternative: store RF energy directly in a high-$Q$
SRF cavity ($Q_0 \sim 5\ee{10}$), then extract via variable
coupler. Key advantages:
\begin{itemize}[nosep]
  \item \textbf{Power gain}: $G = Q_0/Q_{\rm ext}$. At
    $Q_0 = 5\ee{10}$ and $Q_{\rm ext} = 10^5$:
    $G = 5\ee{5}$ --- five orders of magnitude above SLED.
  \item \textbf{Efficiency}: SRF wall-plug to stored energy
    $\sim$50--65\% (W4i14 CWR). SLED/BOC: $\sim$60--70\%.
    Comparable, but P7 stores for $\sim$1~s (allowing CW
    charging from small sources) vs SLED's $\sim\mu$s fill.
  \item \textbf{Facilities that benefit}: Any pulsed linac
    requiring $>$100~MW peak RF: CLIC ($\sim$50~GW total),
    ILC ($\sim$10~MW/cavity $\times$ 16,000 cavities),
    compact medical linacs (FLASH therapy at $>$1~Gy/pulse),
    FCC-ee injector.
\end{itemize}
P7 does \textbf{not} replace klystrons (which are CW/long-pulse
RF sources); it replaces the \emph{pulse compression system}
downstream of klystrons by providing a fundamentally higher
compression ratio. \textbf{Estimated accelerator RF TAM:
\$0.3--1.0~B over the next decade (CLIC R\&D, ILC if approved,
compact linac market).}

\textbf{Finding 4 [HIGH IMPACT --- NEW]: Space-Based DEW ---
P7's 371~kg Conduction-Cooled System Matches SDA Power Density
Requirements.}
DoD is scaling space-based directed energy from 150~kW (current)
to 500~kW (2025--2030) and eventually MW-class. The Space
Development Agency's Tranche~3 tracking layer (\$3.5~B, 72
satellites, launch 2029) will require point defense capability.
Space DEW's critical constraint is power-to-mass ratio:
\begin{itemize}[nosep]
  \item \textbf{Solar arrays}: $\sim$100--300~W/kg. A 500~kW
    system needs 1,700--5,000~kg of solar panels alone,
    \emph{plus} energy storage for pulsed operation.
  \item \textbf{Batteries (Li-ion)}: 150--250~Wh/kg $\approx$
    0.5--0.9~MJ/kg. A 10~MJ magazine: 11--20~kg. But
    discharge rate limited ($<$10C $\Rightarrow$ $>$6~s),
    cycle life $\sim$500--2000 deep cycles.
  \item \textbf{Capacitors}: 0.01--0.05~MJ/kg. A 10~MJ
    magazine: 200--1000~kg. Fast discharge but heavy.
  \item \textbf{P7 SRF (conduction-cooled Nb$_3$Sn at 371~kg,
    W4i14)}: 46~MJ stored $\Rightarrow$ 124~kJ/kg effective
    energy density (\emph{including} cryocooler, structure,
    thermal management). Discharge in $<$1~ms.
    Cycle life: $>10^7$.
\end{itemize}
P7's 124~kJ/kg exceeds capacitors by $\sim$2,500--12,000$\times$
but is below Li-ion on a per-kg basis for small energy stores.
The conduction-cooled architecture (no liquid helium) is uniquely
suited to microgravity. China has reportedly begun development of
advanced satellite power systems for directed-energy applications.
\textbf{P7's advantage emerges at $>$5~MJ stored energy, where
battery cycle life and capacitor mass become prohibitive.
Space DEW power magazine is a capability gap that P7 uniquely
fills at MW-class. Estimated TAM: \$0.5--2~B (2028--2035),
classified procurement likely.}

\textbf{Finding 5 [MODERATE IMPACT --- NEW]: Fusion Pilot Plant
Auxiliary Systems --- P7 as Plasma Heating and Current Drive
Buffer for Pulsed Fusion Concepts Only.}
The DOE Fusion Science and Technology Roadmap (October 2025) targets
an operating fusion pilot plant by mid-2030s but warns of ``critical
gaps'' in auxiliary systems. P7's 1~s storage time and sub-ms
discharge make it unsuitable for steady-state heating (tokamak
ICRF/ECRF, stellarator). However, it is the ideal technology for:
(1)~Z-pinch driver (Zap Energy: 1--10~MJ per pinch at $>$3~Hz);
(2)~IFE laser driver (Focused Energy, Xcimer: multi-MJ at
10--15~Hz); (3)~Tokamak disruption mitigation (1--10~MJ in
$<$10~ms, safety-critical); (4)~Plasma startup for all concepts.
Three of eight DOE Milestone awardees (Focused Energy, Xcimer,
Zap) have pulsed power requirements that P7 directly addresses.
\textbf{Combined fusion auxiliary TAM (pulsed concepts only):
\$1--4~B (2030--2040).}

\end{tcolorbox}

%=============================================================================
\part{Detailed Analysis}
\label{part:detail}
%=============================================================================

%=============================================================================
\section{Finding 1: IFE Driver Integration}
\label{sec:ife}
%=============================================================================

\subsection{NIF Status (as of March 2026)}

The National Ignition Facility has now achieved ignition (fusion energy
exceeding laser energy delivered to target) on 10 separate occasions:
\begin{itemize}[nosep]
  \item December 2022: First ignition, 3.15~MJ yield (gain $\sim$1.5)
  \item April 2025: Record yield 8.6~MJ, gain $>$4.1 (2.08~MJ laser input)
  \item October 2025: 3.5~MJ yield (10th ignition shot)
\end{itemize}

NIF operates at $\sim$1 shot/day (limited by optics damage, target
fabrication). A commercial IFE plant requires 10--15~Hz. This
$\sim$10$^6\times$ rep-rate gap is the central engineering challenge
for IFE.

\subsection{P7 Integration Architecture for IFE}

\begin{finding}[P7 as IFE Pulsed Power Buffer]
\label{find:ife-buffer}
An IFE plant at 10~Hz with 2~MJ/pulse laser energy requires:
\begin{equation}
P_{\rm avg} = \frac{E_{\rm pulse}}{\eta_{\rm laser}} \times f_{\rm rep}
= \frac{2~\text{MJ}}{0.15} \times 10~\text{Hz} = 133~\text{MW (electrical)}.
\end{equation}
The grid supplies 133~MW CW. P7 acts as the intermediate buffer:
\begin{enumerate}[nosep]
  \item Grid $\to$ CW RF source $\to$ SRF cavity charged over
    $\Delta t = 1/f_{\rm rep} = 100$~ms (well within $\tau_{1/e} = 1$~s).
  \item At $t = 0$: coupler switches to $Q_{\rm ext} = 10^4$,
    cavity discharges 13.3~MJ in $\tau_d = Q_{\rm ext}/\omega
    \approx 2.4~\mu$s (at 650~MHz).
  \item CWR converts RF to DC/pulsed-power format for laser
    diode arrays or Z-pinch capacitor bank.
\end{enumerate}

\textbf{Key advantage}: The SRF cavity smooths the 10~Hz pulsed
demand into CW grid draw. Without P7, the grid sees 133~MW
pulsed at 10~Hz with peak-to-average ratio $>100$, requiring
massive capacitor banks ($>$1000~m$^3$ at 0.01--0.1~MJ/m$^3$)
or dedicated pulsed power supplies.

\textbf{With P7}: A single Nb$_3$Sn multi-cell cryomodule
($\sim$2~m$^3$, storing 208~MJ) replaces $>$1000~m$^3$ of
capacitor banks. Volume reduction: $>$500$\times$.
\end{finding}

\subsection{IFE Companies and P7 Relevance}

\begin{center}
\begin{tabular}{lccc}
\toprule
Company & Concept & Pulsed Power Need & P7 Fit \\
\midrule
Focused Energy & Laser IFE & 10 Hz, multi-MJ & \textbf{Excellent} \\
Xcimer Energy & Laser IFE (excimer) & 10 Hz, multi-MJ & \textbf{Excellent} \\
Zap Energy & Z-pinch & $>$3 Hz, 1--10 MJ & \textbf{Excellent} \\
CFS (SPARC/ARC) & Tokamak & CW (disruption: pulsed) & Marginal \\
Tokamak Energy & Sph.\ tokamak & CW heating & Poor \\
Thea Energy & Stellarator & Steady-state & Poor \\
Realta Fusion & Mirror & Quasi-CW & Poor \\
\bottomrule
\end{tabular}
\end{center}

Three of eight DOE Milestone awardees (Focused Energy, Xcimer, Zap)
have pulsed power requirements that P7 directly addresses.

\subsection{IFE Repetition Rate Technology Status}

Key 10~Hz laser driver milestones:
\begin{itemize}[nosep]
  \item HAPLS (LLNL): 30~J / 30~fs / 10~Hz --- world's first 10~Hz PW laser
  \item HALNA~10 (Japan): 10~J $\times$ 10~Hz demonstrated;
    HALNA~100 (100~J $\times$ 10~Hz) under construction
  \item Colorado State University: 10~Hz high-power laser facility
  \item MEC-U (SLAC): 10~Hz pulse delivery
\end{itemize}

All of these require pulsed power delivery at 10~Hz. P7 is
directly applicable as the intermediate energy buffer between
CW grid power and pulsed laser diode arrays.

%=============================================================================
\section{Finding 2: DEW Cost-Per-Shot Analysis}
\label{sec:dew-cost}
%=============================================================================

\subsection{Cost Model}

The cost per shot for a pulsed energy weapon is:
\begin{equation}
C_{\rm shot} = \underbrace{\frac{C_{\rm capital}}{N_{\rm life}}}_{\rm amortization}
+ \underbrace{C_{\rm energy}}_{\rm electricity}
+ \underbrace{C_{\rm maint}/N_{\rm year}}_{\rm maintenance}
\end{equation}

For a 10~MJ magazine (sufficient for a 100~kW HEL firing for 100~ms
or an HPM burst):

\begin{center}
\begin{tabular}{lcccc}
\toprule
Technology & $C_{\rm capital}$ & $N_{\rm life}$ & $C_{\rm energy}$ & $C_{\rm shot}$ \\
\midrule
Capacitor bank & \$500K--2M & $10^4$--$10^5$ & \$0.003 & \$0.50--2.00 \\
Flywheel & \$1--3M & $>10^6$ & \$0.003 & \$0.10--0.30 \\
P7 SRF (Nb) & \$2--5M & $>10^7$ & \$0.003 & \textbf{\$0.02--0.08} \\
P7 SRF (Nb$_3$Sn) & \$3--8M & $>10^7$ & \$0.003 & \textbf{\$0.03--0.12} \\
\bottomrule
\end{tabular}
\end{center}

\textbf{Cryogenic OPEX}: P7's cryogenic operating cost is
$\sim$6~kW (conduction-cooled) $\times$ 8,760~hr/yr $\times$
\$0.10/kWh $= \$5,300$/yr. Amortized over $\sim$10$^5$ shots/yr
(typical DEW training + engagement tempo): \$0.05/shot.

\textbf{Key insight}: P7's advantage comes from superconductor
cycle life ($>10^7$, limited only by cryocooler maintenance)
vs capacitor dielectric degradation ($10^4$--$10^5$ cycles).

\subsection{Volume and Mass Comparison for Mobile Platforms}

\begin{center}
\begin{tabular}{lccc}
\toprule
Technology & 10 MJ Volume & 10 MJ Mass & Discharge Time \\
\midrule
Capacitor bank & 100--1000 m$^3$ & 5,000--50,000 kg & $<$1 ms \\
Flywheel & 0.05--0.2 m$^3$ & 20--100 kg & $>$100 ms \\
P7 SRF (Nb) & 0.015 m$^3$ & $\sim$200 kg (w/ cryo) & $<$1 ms \\
P7 SRF (Nb$_3$Sn) & 0.003 m$^3$ & $\sim$100 kg (w/ cryo) & $<$1 ms \\
\bottomrule
\end{tabular}
\end{center}

P7 uniquely combines the fast discharge of capacitor banks with
the compactness of flywheels, while exceeding both on cycle life.

%=============================================================================
\section{Finding 3: Accelerator RF Pulse Compression}
\label{sec:accel}
%=============================================================================

\subsection{Current State of the Art}

The CLIC pulse compression system (2025 design, Phys.\ Rev.\ Accel.\
Beams 28, 032002) uses SRF storage cavities and correction cavities
to achieve power gain $G = 3.78$. P7 proposes a fundamentally different
approach: \emph{accumulate} RF energy over many fill cycles
($\sim$10$^9$ oscillations at $Q_0 = 5\ee{10}$) and then extract
in a single discharge:
\begin{equation}
G_{\rm P7} = \frac{Q_0}{Q_{\rm ext}} = \frac{P_{\rm peak}}{P_{\rm charge}}.
\end{equation}
At $Q_0 = 5\ee{10}$, $Q_{\rm ext} = 10^5$: $G_{\rm P7} = 5\ee{5}$.

\subsection{Comparison with SLED/BOC}

\begin{center}
\begin{tabular}{lcccc}
\toprule
System & Power Gain & Efficiency & Fill Time & Complexity \\
\midrule
SLED (SLAC) & 2--4$\times$ & 60--65\% & $\sim\mu$s & Low \\
BOC (KEK) & 3--5$\times$ & 55--65\% & $\sim\mu$s & Medium \\
CLIC pulse comp. & 3.78$\times$ & $\sim$70\% & $\sim\mu$s & High \\
\textbf{P7 SRF} & $\mathbf{5\ee{5}\times}$ & 50--65\% & $\sim$1 s & Medium \\
\bottomrule
\end{tabular}
\end{center}

\textbf{Caveat}: P7's 1~s fill time means it cannot follow
pulse-to-pulse variations at $>$1~Hz linac rep rates without
multiple cavities in rotation. For a 5~Hz linac (ILC), 5~P7
cavities in round-robin would each have 1~s to charge.

\subsection{Facilities That Would Benefit}

\begin{enumerate}[nosep]
  \item \textbf{CLIC} (if approved): Replaces entire pulse
    compression chain. Reduces klystron count by $\sim$100$\times$.
  \item \textbf{Compact medical linacs}: FLASH radiotherapy
    requires $>$40~Gy/s. P7 enables single-shot delivery.
  \item \textbf{Industrial electron beam processing}: Pulsed
    sterilization, materials modification at $>$100~kW peak.
  \item \textbf{FCC-ee injector}: Very high peak power
    for booster ring injection.
\end{enumerate}

\textbf{Important clarification}: P7 does \textbf{not} replace
klystrons (which are CW/long-pulse RF sources). It replaces the
\emph{pulse compression system} downstream of klystrons.

%=============================================================================
\section{Finding 4: Space-Based DEW}
\label{sec:space-dew}
%=============================================================================

\subsection{DoD Space DEW Roadmap}

The DoD DE roadmap (May 2024) targets:
\begin{itemize}[nosep]
  \item 2025--2030: Scale from 150~kW to 500~kW
  \item 2030+: MW-class directed energy
  \item SDA Tranche~3: \$3.5~B for 72 tracking satellites (2029 launch)
\end{itemize}

\subsection{P7 Space DEW Architecture}

The W4i14 conduction-cooled Nb$_3$Sn system (371~kg total,
46~MJ stored, 6~kW cryocooler power) provides:
\begin{itemize}[nosep]
  \item 46~MJ / 50~kJ per engagement = 920 engagements
    from full charge (without recharge)
  \item Recharge from solar arrays: 46~MJ / 50~kW solar
    $= 920$~s $\approx$ 15~min for full recharge
  \item Burst mode: 10 engagements in 10~s (each 50~kJ,
    total 500~kJ $\ll$ 46~MJ stored)
\end{itemize}

\textbf{Mass comparison for 500~kJ burst capability}:
\begin{center}
\begin{tabular}{lcc}
\toprule
Technology & Mass (500 kJ) & Volume \\
\midrule
Li-ion battery & 2--4 kg & 0.002--0.004 m$^3$ \\
Supercapacitor & 50--100 kg & 0.05--0.1 m$^3$ \\
Capacitor bank & 10--50 kg & 5--50 m$^3$ \\
P7 SRF (371 kg total) & 371 kg (stores 46 MJ) & 0.07 m$^3$ \\
\bottomrule
\end{tabular}
\end{center}

\textbf{Key insight}: For the 500~kJ burst case, Li-ion batteries
are lighter. P7's advantage emerges at $>$5~MJ stored energy
(where battery cycle life and capacitor mass become prohibitive)
and at MW-class power levels where discharge speed matters.

\textbf{P7 sweet spot for space DEW}: MW-class weapons
requiring $>$10~MJ per engagement, $>10^4$ lifetime shots.

%=============================================================================
\section{Finding 5: Fusion Pilot Plant Auxiliary Systems}
\label{sec:fpp}
%=============================================================================

\subsection{DOE Fusion Roadmap (October 2025)}

The DOE Fusion Science and Technology Roadmap identifies ``critical
gaps'' in technologies needed for a fusion pilot plant by mid-2030s.
The GAO (GAO-25-107037, 2025) found that additional planning is
needed to strengthen DOE's fusion commercialization efforts.

Key milestones:
\begin{itemize}[nosep]
  \item 8 companies in Milestone program (Period~1 completing
    late 2025, Period~2 underway 2026)
  \item At least 3 companies completed $\geq$5 S\&T milestones
    by January 2025
  \item Thea Energy: first DOE-certified preconceptual design
    (January 2026)
\end{itemize}

\subsection{P7 Role in Pulsed Fusion Concepts}

\begin{finding}[P7 Natural Niche: Pulsed Fusion Only]
\label{find:fpp-niche}
P7's 1~s storage time and sub-ms discharge make it unsuitable
for steady-state heating. Ideal applications:
\begin{enumerate}[nosep]
  \item \textbf{Z-pinch driver} (Zap Energy): 1--10~MJ per pinch
    at $>$3~Hz.
  \item \textbf{IFE laser driver} (Focused Energy, Xcimer):
    Multi-MJ per pulse at 10--15~Hz.
  \item \textbf{Tokamak disruption mitigation}: 1--10~MJ
    in $<$10~ms for massive gas injection / shattered pellet
    injection. Rare but safety-critical.
  \item \textbf{Plasma startup} (all concepts): 1--50~MJ, $<$1~s.
\end{enumerate}
\end{finding}

%=============================================================================
\section{Updated TAM Summary}
\label{sec:tam}
%=============================================================================

\begin{center}
\begin{tabular}{lcc}
\toprule
Segment & TAM (2025--2035) & Confidence \\
\midrule
DEW power magazine & \$2.4--5.6~B & High \\
IFE pulsed power & \$2--8~B & Medium (contingent) \\
Accelerator RF pulse compression & \$0.3--1.0~B & Medium \\
Space DEW (MW-class) & \$0.5--2~B & Medium-Low \\
Fusion pilot plant aux.\ (pulsed) & \$1--4~B & Medium (contingent) \\
\midrule
\textbf{Committed (DEW + accel)} & \textbf{\$3.2--7.1~B} & High \\
\textbf{Contingent (IFE + space + fusion)} & \textbf{\$3.0--13.5~B} & Medium \\
\textbf{Total potential} & \textbf{\$6.2--20.6~B} & --- \\
\bottomrule
\end{tabular}
\end{center}

%=============================================================================
\section{Corrections and Consistency Checks}
\label{sec:corrections}
%=============================================================================

\begin{correction}[TAM Range Refinement]
\label{corr:tam}
W4i14 reported P7 TAM as \$4.1--9.7~B (2025). With IFE and space
DEW segments added (both contingent), the potential TAM expands to
\$6.2--20.6~B. The committed TAM (DEW + accelerator) is
\$3.2--7.1~B --- consistent with W4i14 lower bound. No W4i14 or
W4i15 quantities contradicted.
\end{correction}

\begin{correction}[Nb$_3$Sn Pulsed TRL Clarification]
\label{corr:trl}
W4i15 reported Nb$_3$Sn TRL 4--5 based on JLab Gray Enid I.
Clarification: this applies to CW \emph{accelerator} use only.
Nb$_3$Sn for \emph{pulsed energy extraction} remains at TRL~1
(no discharge experiment performed). The W4i14 statement
``Nb$_3$Sn pulsed TRL 1 is single highest risk'' remains correct.
\end{correction}

\begin{correction}[Space Energy Density Units]
\label{corr:units}
The W4i16 Finding~4 executive summary initially stated P7's effective
energy density as 124~MJ/kg. This is incorrect. 46~MJ / 371~kg
= 124~kJ/kg. Corrected in body text. This is still $\sim$140--250$\times$
higher than batteries on a per-kg basis for the full 46~MJ store
(batteries would need $\sim$50,000--90,000~kg for 46~MJ including
thermal management and discharge electronics), but the per-kg
advantage only manifests at large stored energies.
\end{correction}

All other inherited facts confirmed. No W4i15 findings contradicted.

%=============================================================================
\section{Derivation Chains}
\label{sec:chains}
%=============================================================================

\textbf{Finding 1 chain}:
DFD $W$-function $\to$ multi-mode SRF storage $\to$ 46--208~MJ
in $\sim$0.015--0.07~m$^3$ $\to$ sub-ms discharge via variable
$Q_{\rm ext}$ $\to$ CWR 80--90\% RF-to-DC $\to$ IFE driver
requires 10~Hz, multi-MJ $\to$ P7 is only technology combining
energy density + rep-rate + discharge speed.

\textbf{Finding 2 chain}:
Superconductor cycle life $>10^7$ (no dielectric fatigue) $\to$
amortized capital cost $<$\$0.01/shot $\to$ add cryogenic OPEX
\$0.05/shot $\to$ total \$0.02--0.08/shot $\to$ 10--100$\times$
below capacitor banks.

\textbf{Finding 3 chain}:
SRF $Q_0 = 5\ee{10}$ (demonstrated) $\to$ store RF for $\sim$1~s
$\to$ extract at $Q_{\rm ext} = 10^5$ $\to$ power gain
$G = Q_0/Q_{\rm ext} = 5\ee{5}$ $\to$ exceeds SLED by
$\sim$10$^5\times$.

\textbf{Finding 4 chain}:
Conduction-cooled Nb$_3$Sn (W4i14: 371~kg, 46~MJ) $\to$
no liquid He (microgravity compatible) $\to$ 124~kJ/kg
$\to$ advantage at $>$5~MJ stored energy $\to$
MW-class space DEW sweet spot.

\textbf{Finding 5 chain}:
$\tau_{1/e} = 1$~s + sub-ms discharge $\to$ ideal for
$>$1~Hz pulsed fusion (Z-pinch, IFE) $\to$ poor fit for
CW/long-pulse (tokamak heating, stellarator) $\to$ P7
niche = pulsed fusion concepts $\to$ 3/8 DOE Milestone
companies.

\end{document}
